
\section{Introduction}\doHeader{Introduction}

\paragraph{Why algorithms?}
Information surrounds us, but much of it is out of our reach. 
Here's a simple example.
Suppose you want to know the prime factors of the following 186-digit number:
{\footnotesize
11679847981112819759721399310592745791658017001955007325132913837831330495881519756453703742878526148\textbackslash\\84146888067442512219413748768010657572575384986457405973985247465176041951676954461208131403777}. 
Like every positive integer, this one has a unique set of prime factors.
In principle you could find the factors by exhaustive search---after all, there are only finitely many possibilities.
But there are about $10^{93}$ possible factors to check,
so that would take too long, even for a computer.
In fact, no algorithm is known that can factor large integers
in any reasonable amount of time
(e.g., the lifetime of the universe).\footnote
{No such algorithm is known---indeed, standard encryption schemes rely on this fact.
  But nobody knows for sure that there isn't one.
  (We do know that, if a so-called quantum computer can be built,
  it would be able to factor efficiently.)}
The information is (in some sense) ``there'', but it is out of reach,
because we don't know an efficient algorithm to compute it.

This kind of barrier to knowledge is common.  For example, Google has a database with details of all the roads in the United States.  From that data, Google Maps can quickly compute, say, a shortest possible route from New York City to San Francisco.  But what about a shortest possible route that \emph{visits all 50 U.S. capitols in any order}?  Nobody knows how to find such a route with a reasonable amount of computation time.

Whether or not we can learn a desired piece of information often depends on whether there is an \emph{efficient algorithm to solve the underlying computational problem}.  Further, for many important and easily stated problems, whether the problem has an efficient algorithm \emph{is not currently known}.  New algorithms can lead to advances in science, medicine, biology, business, and so on.  % This is a fundamental scientific frontier, one that is accessible to you sitting at a desk and working only with paper and pencil.

\paragraph{Why proofs?}
Algorithms are important, but can be difficult to reason about.
For example, consider the following simple algorithm:

\begin{framed}
  \vspace*{-1ex}
  \begin{steps}
  \item[\bf input:] integer $N\ge 1$
    \step while $N > 1$:
    \begin{block}
      \step if $N$ is even:  $N\leftarrow N/2$
      \step else: $N\leftarrow 3*N + 1$
    \end{block}
    \step print ``done''
  \end{steps}
  \vspace*{-1ex}
\end{framed}

In 1937, L.~Collatz conjectured that, for every input, this algorithm eventually halts.
But even after years of study, nobody knows for sure!\,\footnote
{See \url{https://en.wikipedia.org/wiki/Collatz_conjecture}.}
Many simple algorithms are hard to reason about.
So, when searching for an algorithm to solve a particular problem,
we limit our search to algorithms that
we \emph{do} know some way to reason about---that is, algorithms that we can prove do work.
For this reason, \emph{proofs drive algorithm design.}
(More on this later.)
% The mental process of searching for an algorithm is intertwined
% with the process of searching for a proof.
% So, if you want to be good at designing algorithms,
% you must be able to do proofs!
% \,\footnote
% {Of course, correctness is just one of many considerations in algorithm design.
% Other factors such as efficiency and simplicity are also important.  They are just
% not our focus here.}

\paragraph{Do your own pull-ups!}
Textbooks and lectures generally present algorithms that are already very well studied.
The author and reader (or instructor and student) generally \emph{already know} that the algorithms are correct.
So textbooks and lectures focus on conveying the main underlying ideas.
Listening to such a lecture or reading such a text,
it is not hard to get an intuitive sense of the algorithm and why it works.
And in many algorithms courses, this is all you would do.

But this experience does not teach you how to design your own algorithms.
To find an algorithm that works, you typically must search through a variety of ideas.
Ideas that seem promising often fail, sometimes for deviously subtle reasons.
To know that an idea works, you must also find an airtight line of reasoning as to why.
That is, the search for an algorithm proceeds hand-in-hand along with the search for a proof that it works.
Doing this is much harder than following a proof in a lecture or textbook.
The difference is like the difference between watching someone else do pull-ups,
versus doing them yourself.
If you want to build the strength to do pull-ups, it doesn't help to just watch someone else do them.
You have to do them yourself!
So, what does this entail?

\paragraph{Mind the gap!  (Check the details.)}
When you are developing a new proof, you are usually not quite sure it is right.
With experience, you will learn that your intuition has blind spots. 
Because of this, you have to check the details carefully by writing out your reasoning.
Any gap in your \emph{explanation} can hide a gap or error in your \emph{reasoning}.
So, when developing a new proof,
it is important to explain (and check!) your reasoning step by step,
in more detail than you see in a typical textbook or lecture.
Also, the more clearly you communicate your reasoning,
the more a reader can understand about the line of reasoning you have in mind,
so the more useful their feedback can be.

\paragraph{Exploring the logical consequences of an assumption.}
Suppose we don't know that a given statement is true.
We can still assume for the moment that the statement is true, and explore the logical consequences if it is.
This kind of imaginative thinking can help figure out what is actually true.
It is so common that keeping track of assumptions
is one of the main organizational challenges in presenting a proof.
Conventionally, this is done via four standard proof patterns,
\emph{proof by contradiction}, reasoning by \emph{cases},
and when proving \emph{if then} and \emph{for all} statements.
When creating complex proofs, it is useful to know these forms
like the back of your hand, so you can keep track of your assumptions
and organize your reasoning clearly.

\paragraph{These notes: \emph{long-form} proofs.}
To learn to develop proofs with enough detail,
and to learn the four standard forms for managing assumptions,
we'll use \emph{long-form proofs}.
A long-form proof is given as a sequence of numbered steps,
where each step is small, has a clear meaning, and follows easily from previous steps.
The numbering allows each step to be identified unambiguously.
(This is useful during the proof to make it easy to tell the reader which previous steps a particular step follows from.
It also makes communicating about the proof, including giving feedback on it, easier.)

In long-form proofs, assumptions are explicitly managed by using what we call \emph{blocks}.  
Each block introduces an assumption, which then holds for the scope of the block.
Blocks nest like code blocks in computer programs.
There are four block types---one for each of the four standard proof forms mentioned above. 
Section~\ref{proofs: sec: long form} describes the long-form proof format in detail.

\paragraph{Basic terms: \emph{problem}, \emph{algorithm}, \emph{proof}}
To design algorithms, and to reason and communicate clearly about them,
it is crucial that the language we use is clear,
with a precise and unambiguous meaning.
To this end, we use the language of mathematics.

A computational \textbf{\emph{problem}} is defined by
specifying what the possible inputs are,
and what the output should be for each input.
(This means that defining a problem is just defining a function.)
We cast our problems in terms of standard mathematical objects
such as numbers, strings, trees, graphs, sequence, sets and so on.

As described previously in \LN{prelim},
to define an \textbf{\emph{algorithm}}, one specifies exactly how,
given any input, the algorithm processes the input to produce the output.
The definition should be mathematically precise,
and written to be clear and accessible to any competent reader.
For this reason, use words, not code.
Any competent programmer should be able to translate the definition into code.
An algorithm is \textbf{\emph{correct}} (for a given problem) if,
for every input, the algorithm produces a correct output.


A \textbf{\emph{proof}} of a fact is a clear explanation, typically in writing,
of a complete and correct line of reasoning that establishes that the fact must be true.
Typically, the proof is preceded by a clear statement of the fact being proven,
e.g.~stated as a theorem or lemma.
%
Before you try to write a proof for a fact,
you should already have in mind a line of reasoning as to why the fact must be true.\footnote
{But of course it is also useful
to write partial proofs to guide your thinking.}
The reasoning must be complete (covering all cases) and without logical gaps.
The goal of the written proof is to communicate your line of reasoning to the reader.
The fundamental guiding principle is this:
\emph{from what you write,
  the reader should be able to easily reconstruct your reasoning in detail,
  and to easily verify that it is correct.}
% We say much more later about how to write proofs.
% We use a particular form that we call \emph{long form}.


%%% Local Variables:
%%% mode: latex
%%% TeX-master: "long_form_proofs"
%%% End:
